\begin{pastexam}[label={px:ust-115-ee-ds-37}]{臺聯大電機類 115 學年度 第 37 題}
（單選；2 分）

A magic square is an $n\times n$ matrix of the integers from 1 to $n^2$ such that the sum of each
row and column and the two major diagonals is the same. The following figure shows a magic square
for the case $n=5$ where the common sum is 65.

\begin{center}
\begin{tabular}{|c|c|c|c|c|}\hline
15 & 8 & 1 & 24 & 17\\\hline
16 & 14 & 7 & 5 & 23\\\hline
22 & 20 & 13 & 6 & 4\\\hline
3 & 21 & 19 & 12 & 10\\\hline
9 & 2 & 25 & 18 & 11\\\hline
\end{tabular}
\end{center}

\textbf{Coxeter} has given the following rule for generating a magic square when $n$ is odd:\\
1. Put a \textbf{one} (\texttt{// 1.}) in the middle box of the top row.\\
2. Go up (\texttt{// 2.1}) and left (\texttt{// 2.2}) assigning numbers in increasing order to
empty boxes. If your move causes you to jump off the square (that is, you go beyond the square's
boundaries), figure out where you would be if you landed on a box on the opposite side of the
square. Continue with this box.\\
3. If a box is occupied (\texttt{// 3.1}), go down (\texttt{// 3.2}) instead of up and continue.
\begin{minted}[linenos, numbersep=1.5mm, xleftmargin=2em, fontsize=\footnotesize]{c}
square[0][(n-1)/2]=1;            // 1. Top row [0]; middle box [(n-1)/2]
/* i and j are current position */
i=0;
j=(n-1)/2;
for (count=2;count<=n*n;count++) {
  row=(i-1<0)?(n-1):(i-1);         // 2.1 up
  column=(j-1<0)?(n-1):(j-1);      // 2.2 left
  if(square[row][column])          // 3.1 occupied
    i=(++i)%n;                     // 3.2 down
  else{
    i=_______ ;                    // 3.3 _______
    j=(j-1<0)?(n-1):--j;
  }
  square[i][j]=count;
}
\end{minted}
Fill the blank \texttt{i=\_\_\_\_\_\_\_;} with correct code at the location annotated by
\texttt{//3.3}.

(A)~\texttt{n} \qquad (B)~\texttt{row} \qquad (C)~\texttt{column} \qquad (D)~\texttt{count}
\qquad (E)~\texttt{i\%n}
\end{pastexam}

\begin{pastsoln}
答案是 (B)。

每一輪先算出「往上、往左一格」的位置 \texttt{(row, column)}，超出邊界就繞到對面（第 6、7 行）。
如果那一格已經有數（第 8 行），就改成往下一格：列加一、行不變（第 9 行）。 否則（第 10 行的
\texttt{else}）就移到那一格，列要變成 \texttt{row}、行變成 \texttt{column}。 第 12 行的
\texttt{j=(j-1<0)?(n-1):--j;} 算的正是 \texttt{column}，所以空格要填 \texttt{row}，註解 3.3 是
``up''。

用 $n = 5$ 追蹤前幾步，和題目的方陣對照（列、行從 $0$ 數）：
\begin{center}
\begin{tabular}{cccl}
\toprule
\texttt{count} & 上一格 & \texttt{(row, column)} & 結果 \\
\midrule
1 & & & $(0, 2)$，第 1 行 \\
2 & $(0, 2)$ & $(4, 1)$，空 & 放在 $(4, 1)$ \\
3 & $(4, 1)$ & $(3, 0)$，空 & 放在 $(3, 0)$ \\
4 & $(3, 0)$ & $(2, 4)$，空 & 放在 $(2, 4)$ \\
5 & $(2, 4)$ & $(1, 3)$，空 & 放在 $(1, 3)$ \\
6 & $(1, 3)$ & $(0, 2)$，有 $1$ & 往下，放在 $(2, 3)$ \\
\bottomrule
\end{tabular}
\end{center}
題目的方陣中 $2$ 在第 5 列第 2 行、$3$ 在第 4 列第 1 行、$4$ 在第 3 列第 5 行、$5$ 在第 2 列第 4 行、
$6$ 在第 3 列第 4 行，全部吻合；填 \texttt{row} 一路做完，得到的正是題目那個方陣。

逐項檢查：
\begin{itemize}
\item (A) \texttt{n}：索引超出 $0..n-1$。 錯。
\item (B) \texttt{row}：對。
\item (C) \texttt{column}：把行號當成列號，位置錯。 錯。
\item (D) \texttt{count}：是要填的數字，不是位置；$\texttt{count} \ge 2$ 很快就超出範圍。 錯。
\item (E) \texttt{i\%n}：等於 \texttt{i}，列不動，只往左移，不是「往上往左」。 錯。
\end{itemize}
（附註：第 9 行 \texttt{i=(++i)\%n;} 在同一個運算式裡改了 \texttt{i} 兩次，嚴格說在 C 裡是未定義
行為，應寫成 \texttt{i=(i+1)\%n;}。 題目的意思是往下一格，不影響本題。）
\end{pastsoln}
